\section{Setting: which channel does $\theta=2/9$ live in?}
The Koide angle $\theta=2/9$ is confirmed empirically to seven figures but is not derived from
the geometry. It is itself not an ansatz but a \emph{found} quantity: only the Hilbert-space
translation of FFGFT (Docs.~230/231/232) and its continuation to the masses (Doc.~282) made
visible that the three lepton masses are the spectrum of a hermitian Z$_3$ circulant on the
$\mathbb{C}^3$ fibre, in whose eigenvalues $\theta=2/9$ appears as the result of the
diagonalisation. It is precisely this circulant that supplies the amplitudes $a_k$ below. For the
empirical placement of the angle -- $2/9$ as the rational address of a value fixed via one
reference point, not as a pole-mass exactness claim (P42) -- see Doc.~292; for the dynamical
locus sought here this fine distinction is immaterial, since the orbit resolution lies far above
the $10^{-7}$ level. The preceding analysis (Doc.~290) separates, at the mass operator, a
\emph{magnitude channel} (the eigenvalues, the masses) from a \emph{phase channel} (mixing,
all-pass, holonomy), and shows that a magnitude spectrum fixes the phase only up to an all-pass
factor. The question here is accordingly: in which channel does $\theta=2/9$ sit -- and what
kind of quantity is it, compared with a forced anchor such as the speed of light $c$? The
supporting computations are in \texttt{allpass\_theta.py}, \texttt{allpass\_carrier\_selection.py}
and \texttt{forced\_vs\_contingent.py} (numpy-only, seed 20780458).

\section{The magnitude side is counter-blind}
With the signed amplitudes $a_k(\theta)=1+\sqrt2\cos(\theta+2\pi k/3)$ ($k=0,1,2$), the
symmetric Koide invariant $Q=(\sum_k a_k^2)/(\sum_k a_k)^2$ is exact and
\emph{$\theta$-independent}. The sum $S=\sum_k a_k=3+\sqrt2\sum_k\cos(\theta+2\pi k/3)=3$ (three
cosines spaced by $120^\circ$ sum to zero), and the square sum
$T=\sum_k a_k^2=3+0+2\cdot\tfrac32=6$ (with $\sum_k\cos^2=\tfrac32$). Hence
\[
Q(\theta)=\frac{T}{S^2}=\frac{6}{9}=\frac{2}{3}\qquad\text{for every }\theta.
\]
$Q$ ranges over $[1/3,1]$ for positive masses; the value $2/9=0.2222$ lies \emph{below} that
range. So $2/9$ is no boundary, no extremum, and no fixed point of the symmetric functional --
in the magnitude channel $2/9$ is entirely invisible. This is the algebraic side of the
counter-blindness of the $Z_3$ core: the orbit members $\{1/9,2/9,4/9\}$ are indistinguishable
to any symmetric, counting quantity.

\section{The all-pass: three channels}
The locus of $\theta$ can be made visible at the all-pass. A $Z_3$-symmetric Blaschke product
$B_\theta(z)=\prod_k(z^{-1}-\bar a_k)/(1-a_k z^{-1})$ with poles $a_k=r\,e^{i(\theta+2\pi k/3)}$
(here $r=0.5$) splits the phase information into three quantities, two of which are blind to
$\theta$ and only one of which carries it. Numerically (for the orbit angles
$\theta\in\{1/9,2/9,4/9\}$):
\begin{itemize}
\item \textbf{Magnitude} $|B_\theta|=1$ (to $\sim10^{-15}$) for all $\theta$ -- the magnitude
  channel is blind.
\item \textbf{Winding number} $=3$ exactly for all $\theta$ -- the \emph{flat-topological}
  integer is blind. This is the numerical face of the statement that $2/(9\pi)$ is irrational
  and $\theta=2/9$ is therefore not a flat topological invariant: the integer winding does not
  separate $2/9$ from $1/9,4/9$.
\item \textbf{Continuous phase profile} $\arg B_\theta(\omega)$ -- different for $1/9,2/9,4/9$
  (at the reference point $\omega=0$: $0.093$, $0.171$, $0.249$; the group delay separates them
  too). Only the \emph{continuous, dynamical} phase profile carries $\theta$.
\end{itemize}
An energy/magnitude functional such as $\int|B|^2\,d\omega/2\pi$ is likewise identical for all
three. The distinction of $2/9$ from $1/9,4/9$ thus lives neither in the magnitude, nor in the
integer winding, nor in any symmetric functional -- only in the dynamical phase profile.

\section{An illustrative selecting mechanism}
What a selector would have to do is shown by a dynamical phase reference. The three orbit
members have all-pass phases $\psi=(0.093,\,0.171,\,0.249)$. A carrier that accumulates a
holonomy $\Phi$ and phase-locks to the nearest orbit phase selects exactly one representative:
$2/9$ is selected precisely when $\Phi$ lies in the $2/9$ basin $[0.132,\,0.210]$ (width
$0.078$). A \emph{random} holonomy over the orbit range selects $2/9$ with probability
$\approx1/3$ -- nothing prefers the value $2/9$ without a reason.

The form of this mechanism is that of the open problem: the selection is \emph{contingent} on
an external dynamical phase (a holonomy); $2/9$ has its own basin like any member, and neither
magnitude nor winding privileges it. A dynamical selector would have to supply that holonomy
\emph{forced}, from an internal reason of the supplying dynamics, not by a hand-chosen
relabeling. This is the same
content as the breaking condition $(\delta,\chi)\approx(0.24,\pi/2)$: $\chi$ \emph{is} the
holonomy phase, $\delta$ its free depth. (An illustration of the mechanism shape; not a
derivation of the value $2/9$, and not a model of any specific carrier.)

\section{Forced or contingent: $\theta^*(\delta)$ and the photon contrast}
Whether $2/9$ is \emph{forced} like $c$ is decided by the behaviour under variation. The
$\delta=0$ part of the breaking functional has the closed form (confirmed to $10^{-15}$)
\[
E_3(\theta,0)=-\tfrac12+\tfrac{1}{\sqrt2}\cos 3\theta,\qquad
\frac{dE_3}{d\theta}\Big|_{0}=-\frac{3}{\sqrt2}\sin 3\theta,
\]
with extrema at the $Z_3$ points ($\sin3\theta=0$). Coupling the second triplet $3'$ at
strength $\delta$ and phase $\chi$, the inner extrema move away from the $Z_3$ point (there are
two; the lower, orbit-near branch $\theta^*$ is the one tracked). On the branch $\chi=\pi/2$ the
shift is smooth and monotonic (numerical,
\texttt{forced\_vs\_contingent.py}): $\theta^*$ rises from $0.099$ ($\delta=0.10$) through
$0.223$ ($\delta=0.24$) to $0.501$ ($\delta=1.0$); $d\theta^*/d\delta$ is positive throughout
(median over the $\delta$ grid $0.31$, locally at $\delta^*\approx0.24$ about $0.80$). The
reachable range $[0.05,\,0.52]$ contains $2/9$ in its interior, unremarkable. The $2/9$ basin in
$\theta$ space -- nearest orbit member, $[1/6,\,1/3]$, width $1/6$ -- corresponds via the local
slope to a $\delta$-window of $\Delta\delta\approx0.167/0.80\approx0.21$.
$\delta$ is soft, not fine-tuned; $2/9$ is dragged in, it sits at no structure.

The speed of light is of a different type. In $ds^2=-c^2dt^2+dx^2+dy^2+dz^2$, $c$ is a metric
coefficient, not a parameter of a functional; $ds^2=0$ is the light cone, and massless means
$v=c$ exactly -- the saturation of the inequality $v\le c$. There is no $c^*(\lambda)$;
$dc/d\lambda\equiv0$, because $c$ is a constant of the geometry, not an extremum. This reflects
the physical origin: the photon's masslessness is \emph{gauge-symmetry protected} ($m=0$
exactly, hence $v=c$), a \emph{boundary}/fixed-point value. $2/9$, by contrast, is the
irreducibly $Z_3$-\emph{breaking} angle in the \emph{interior} of the orbit -- symmetry-breaking,
not protected, and therefore contingent.

\section{Self-adjoint and skew-adjoint holonomies}
If the selector is a holonomy phase $\chi$, two sectors must be distinguished. In the breaking
term $\delta\cos(2\varphi_k+\chi)$ ($\varphi_k=\theta+2\pi k/3$), $\chi$ separates an
\emph{even} from an \emph{odd} coupling: $\chi=0$ (and $\chi=\pi$) gives $\cos(2\varphi_k)$, an
even, \emph{self-adjoint}, real coupling; $\chi=\pi/2$ gives
$\cos(2\varphi_k+\pi/2)=-\sin(2\varphi_k)$, an odd, \emph{skew-adjoint}, imaginary coupling.

The two sectors have different natural values. The \emph{self-adjoint} holonomies -- flat $Z_3$
Wilson lines ($\chi=2\pi n/3$) times the two spin structures per cycle (offset $0$ or $\pi$) --
are the multiples of $\pi/3$, $\chi\in\{0,\pi/3,2\pi/3,\pi,4\pi/3,5\pi/3\}$; $\pi/2$ is none of
them. The \emph{skew-adjoint} sector, by contrast, has its natural value at $\chi=\pi/2$: the
$i$-phase, $e^{i\pi/2}=i=\sqrt{-1}$, against the real $e^{i\pi}=-1$. The transition from a real
to a skew-adjoint coupling \emph{halves} the phase $\pi\to\pi/2$ -- it takes the square root of
$-1$.

What matters is which sector supplies the selector. The elimination chain (Doc.~282) has ruled
out every static, self-adjoint candidate; the only survivor is the dynamical \emph{skew-adjoint}
phase. The type-consistent value is therefore $\chi=\pi/2$ -- the $i$-phase of the skew-adjoint
sector -- not a multiple of $\pi/3$. And $\chi=\pi/2$ reaches $\theta^*=2/9$ at
$\delta\approx0.24$ (numerically, \texttt{holonomie\_landschaft.py}). The skew-adjoint type by
itself admits both $i$-values, $\chi=\pi/2$ and $\chi=3\pi/2$ ($e^{\pm i\pi/2}=\pm i$); the
choice of branch is an orientation choice (sense of traversal of the cycle). The two branches
are, however, not equivalent: the mirror branch $\chi=3\pi/2$ has no inner extremum near $2/9$
at $\delta\approx0.24$ (extrema at $0.83$ and $1.87$) -- only the $+i$ branch reaches the orbit
region.

$\chi=\pi/2$ is thus not an arbitrary or topologically homeless value, but the skew-adjoint
$i$-value required by the elimination chain -- with no assumed $Z_4$ and no numerical
coincidence. Precisely, \enquote{skew-adjoint} here refers to the \emph{generator}, not the
finite holonomy: the latter is a unitary group element $\exp(K)$ with skew-adjoint $K$
($K^\dagger=-K$); the skew-adjoint character sits in the generating coupling, not in the finite
holonomy matrix itself. The depth $\delta$ nonetheless remains \emph{free}: the sector fixes the
phase $\chi$, not the depth. The open point is thereby sharpened -- not an arbitrary angle, but
whether a skew-adjoint dynamics carries precisely $\chi=\pi/2$ and thereby fixes the depth
$\delta$ for an internal reason \emph{of that dynamics} (recursion-internally, by Doc.~275, no
value can be distinguished).

\section{The depth $\delta$ as the second harmonic}
The mass amplitude is a Fourier series in the fibre angle,
\[
a(\varphi)=\underbrace{1}_{\text{mean}}+\underbrace{\sqrt2\cos\varphi}_{\text{1st harmonic}}
+\underbrace{\delta\cos(2\varphi+\chi)}_{\text{2nd harmonic}},\qquad \varphi=\theta+2\pi k/3.
\]
The first harmonic, with real amplitude $c_1=\sqrt2$, fixes the Koide invariant $Q=2/3$ (it is
the magnitude $r=\sqrt2$). The second harmonic has the complex amplitude $c_2=\delta\,e^{i\chi}$;
its \emph{magnitude} is $\delta=|c_2|$ and its \emph{phase} is $\chi=\arg c_2$. For $\chi=\pi/2$
the coefficient $c_2=i\delta$ is purely imaginary -- the second harmonic is a pure sine, the
skew-adjoint coupling of the previous section.

So $\delta$ and $\chi$ are the magnitude and phase of the \emph{same} quantity $c_2$, but they
belong to different sectors: $\chi=\arg c_2$ lives in the phase sector (fixed to $\pi/2$ by the
elimination chain), $\delta=|c_2|$ in the \emph{magnitude} sector. The depth $\delta$ is thus not
a phase but an amplitude degree of freedom -- the amplitude of the first overtone, physically the
strength with which the second triplet $3'$ (the doubling $6=3\oplus3'$) mixes into the primary
$3$.

This is the same sector that already fixes $|c_1|=\sqrt2$. As the first harmonic amplitude
follows from $Q=2/3$, the second, $|c_2|=\delta$, would have to follow from the overtone content
of the mass-generating fractal waveform (the structure $D_f=3-\xi$, the $3\oplus3'$ doubling).
The ratio $\delta/\sqrt2\approx0.17$ -- the second harmonic is about $17\,\%$ of the first -- is a
pure \emph{shape} quantity, fixed by the waveform, not by the phase. $\delta$ is therefore not a
free tuning parameter but a not-yet-computed magnitude degree of freedom.

The question of $2/9$ thereby separates cleanly into two sectors: the \emph{phase} $\chi=\pi/2$
is skew-adjoint and fixed; the \emph{magnitude} $\delta$ is the second harmonic amplitude and
belongs -- open -- to the same magnitude sector as $\sqrt2$.

\section{$\delta$ as a golden-damped amplitude}
The overtone content can be sharpened: the second harmonic amplitude is a \emph{golden
damping}. The value that gives $\theta^*=2/9$ at $\chi=\pi/2$ is $\delta^*=0.2389$; the golden
attractor $1/\varphi^3=0.2361$ sits just below it. Since $1/\varphi^3=\varphi^{-3}$ is a
\emph{negative} golden power (a decay), $\delta$ is a damping, not an amplification --
$\delta^*<1$ is unambiguously in the damping regime.

The exponent $3$ is not free but $=D$. A single, decompactified $3'$ recursion would give
$1/\varphi^1=0.618$ ($159\,\%$ off); only damping over $D=3$ dimensions (the fractal/spatial
dimension $D_f=3-\xi$, the \emph{full} compact $T^4$) gives $1/\varphi^3$. The matching
exponent $3$ is thus itself an argument for the full compact $T^4$ and against decompactifying
to a single recursion.

The $\xi$-damping recursion $r(k{+}1)=r(k)(1-\xi)$, $r(0)=D_f/3$ (Doc.~275) passes $1/\varphi$
at depth $k^*=\ln\varphi/\xi\approx3609$ and $1/\varphi^3$ at $3\cdot3609\approx10827$ (to
$0.011\,\%$). $\delta^*$ is reached about $89$ \emph{units} earlier -- meaning a difference of
two real-valued log-depths in the decompactified, continuous scale flow ($1/\varphi^3$ at
$10826.2$, $\delta^*$ at $10737.7$; difference $88.5$), not an integer step count. So $\delta^*$
is a $1/\varphi^3$ under-damped by about $89$ log-units, an excess of $+1.187\,\%$ above the
attractor -- both statements are, via $\ln(\delta^*\varphi^{3})/\xi$, one and the same number,
not two independent findings --, consistent with an approach from above.

Whether this excess is the fractal mass correction $K_{\text{frak}}=1-100\xi$ (Doc.~133)
\emph{cannot be decided} at the present level of precision. Its precise form
$(D_f^{\text{eff}}/3)^{D_f^{\text{eff}}/2}$ with $D_f^{\text{eff}}\approx2.973$ gives a factor
$\approx100$ ($1.333\,\%$), while the measured excess is $\approx89\,\xi$ ($1.187\,\%$) -- a gap of
only $0.147\,\%$. Because $\delta^*$ is tied through $\theta^*=2/9$ to the $\tau$-limited angle
determination ($\mathrm{d}\theta^*/\mathrm{d}\delta\approx0.31$), it itself carries a tolerance of
$\approx0.135\,\%$; the two candidates $89\xi$ and $100\xi$ therefore lie only about $1\sigma$ apart
and are \emph{indistinguishable} at current precision. Rejecting $K_{\text{frak}}$ at the $0.15\,\%$
level would be sharper than $\delta$'s own determination tolerance and is thus inadmissible --
exactly parallel to the $N_0$ assessment (Doc.~292: $100/\varphi^2$ sits at $1.6\sigma$ and is
\emph{admissible}, not refuted). $K_{\text{frak}}$ is moreover structurally motivated: the second
harmonic and $K_{\text{frak}}$ spring from the \emph{same} fractal correction $g=\eta(1+\xi f)$
(Doc.~286), of which $K_{\text{frak}}$ is the RG level (Doc.~133). $K_{\text{frak}}$ is thus an
\emph{admissible, structurally motivated candidate} for the excess, not excluded.

Two limits remain. First, admissible is not derived -- the tolerance removes the false rejection but
does not single out $K_{\text{frak}}$ within the band (which holds $89$, $100$ and everything
between); the distinction over an arbitrary band value rests solely on the structural argument, not
the number. Second, the shared $g=\eta(1+\xi f)$ origin is shown in Doc.~286 for the \emph{phase}
channel (holonomy); the transfer to the \emph{magnitude} of the second harmonic ($\delta$) is
plausible but not yet computed. Independently of this, Doc.~275 does not fix the recursion depth in
any case: there $1/\varphi^3$ is a \emph{transit point}, not a fixed point (the only fixed point of
the damping is $0$), and for \emph{every} target $c\in(0,1)$ a depth $k^*(c)$ exists with identical
parameter status. That $\delta^*$ lies near $1/\varphi^3$ is therefore \emph{not} a closing of the
recursion; the fine value remains an admissible but underived candidate.

So the magnitude sector is clarified as far as the structure carries: $\delta$ is the
golden-damped second harmonic, $1/\varphi^3$ over $D=3$ (full $T^4$), with a small excess from
above. Form, order, exponent and direction are fixed; the exact $\delta$ fine value -- the excess
over $1/\varphi^3$ -- is \emph{compatible} with $K_{\text{frak}}=1-100\xi$ within the $\theta$
tolerance and structurally motivated (Doc.~286, shared $g=\eta(1+\xi f)$ origin), but not derived as
a magnitude value and remains open in that sense (\texttt{delta\_golden\_daempfung.py}).

\section{What the non-closure means in practice: an open coupling channel}
The picture so far separates two sectors with opposite behaviour. The mass sector is
\emph{closed}: the recursion $R(\Phi^*)=\Phi^*$ (Doc.~203) has a non-trivial fixed point, and
because the lepton masses sit on it, they and the Koide value $Q=2/3$ are exact and rigid.\footnote{\enquote{Exact} here refers to the geometric fixed-point level: the dimensionless invariant $Q=2/3$ is rigid and is confirmed by the data to $-0.9\,\sigma$ (Doc.~292). The \emph{absolute} masses, or the number pair $(\sqrt2,\ 2/9)$, must be distinguished from this: on the measured pole-mass level the pair is not jointly exact (Doc.~292, $\sim450\,\sigma$ in the $m_\mu/m_e$ channel), which is why $2/9$ is treated per P42 as the rational address of an angle fixed via one reference point, not as a pole-mass exactness claim. For the dynamical locus discussed here this is immaterial -- the orbit resolution ($\sim0.1$) lies far above the $10^{-7}$ difference.} There
is nothing to adjust there -- no free parameter, no point of entry for an external influence.

The room sits solely in the open sector, the phase. The $\xi$-damping recursion never closes
(transit point, only fixed point $0$, Doc.~275); in feedback language it is an open, damped loop
with gain below one, not a rigidly locked fixed point. Such an open, damped loop can -- unlike a
closed fixed point -- be influenced from outside: by coupling to a second oscillator, by injection
or mode-locking (Doc.~286). In practice the non-closure therefore means that a channel exists
through which a second, dynamical system can lock the open phase. The closed mass sector offers no
such channel.

The channel is, however, narrow and structured, not a free parameter. The capture range of such a
coupling is of order $\xi$ (Doc.~286) -- tiny; the $2/9$ Arnold tongue is about twelve times
narrower than $1/3$. The coupling of adjacent scale levels in the continuous flow is real but of
order $10^{-13}$ in magnitude and does not destabilise the closed predictions. And the effect is
directed: it runs from the dynamics into the phase, not into the radial mass sector, which stays
rigid (the radial/angular split of Doc.~282).

The one place this channel would be used is a dynamical skew-adjoint generator outside the static
FFGFT geometry -- the continuation of the elimination chain (Doc.~283/286). If such a generator
locks the open phase precisely to $\chi=\pi/2$ and thereby $\theta=2/9$, then $2/9$ would acquire
the dynamical, protected status the following conclusion describes. This is explicitly a candidate,
not an attained result: FFGFT's own recursion $R$ is purely real and carries no such generator;
whether an external dynamics supplies one remains open.

The decisive point is the boundary of this room. \enquote{Room for mutual influence} does not mean
the theory can be bent at will. The room is bounded: the radial sector stays rigid, the phase is
tightly constrained (skew-adjoint, $\chi=\pi/2$, transcendental in $2\pi$, $\xi$-narrow capture
range), and the recursion depth fixes no value (Doc.~275). It is structured room, not a free
parameter -- and it is exactly this boundary that distinguishes a channel for genuine coupling from
an after-the-fact numerical adjustment.

This channel is also the reason the fundamental sector does not collapse to a rigid skeleton. A
fully closed model -- every value a forced fixed point -- would be over-determined: no contingent
selection within a $Z_3$ orbit, no flavour mixing, no physical (non-removable) phase. Mixing in
particular is phase, not magnitude (Doc.~289: charged leptons magnitude-dominated with small
mixing, neutrinos phase-dominated with large mixing, the pure all-pass sector). The sharpest case
is the structure of $CP$ violation: it requires an irreducible complex phase -- a purely real,
self-adjoint structure is $CP$-conserving -- hence exactly the type of the skew-adjoint channel
$\chi=\pi/2$. The open phase channel is thus structurally of the kind that can carry $CP$
violation; a fully self-adjoint FFGFT would be $CP$-even and non-mixing. The scope stays narrow:
the channel carries mixing, $CP$ and contingency in the \emph{fundamental sector} -- not the
macroscopic diversity (chemistry, structure formation), which has its own independent sources.

\section{Conclusion: the open dynamical question}
Statically the situation is clear: $2/9$ is invisible in the symmetric $Q$, blind to magnitude
and to the integer winding, and an interior extremum dragged in by a free parameter
($d\theta^*/d\delta\neq0$). It is \emph{not} forced like $c$. Both quantities share the same
flat $T^4$ arena -- $c$ as a signature boundary, $2/9$ as an interior phase degree of freedom --
but they are of different types.

For $2/9$ to acquire a $c$-like, forced status, it would need a \emph{dynamical} fixed point,
as gauge symmetry forces $m=0$ (and hence $v=c$). The preceding sections localise this fixed
point: the \emph{phase} $\chi=\pi/2$ is the skew-adjoint $i$-value required by the elimination
chain -- the selector is type-fixed; what remains open is only the \emph{magnitude} $\delta$
(the golden-damped overtone amplitude). Its fine value is, moreover, not \enquote{not yet}
derived but \emph{provably not} derivable from the recursion alone: the non-closure is a theorem
(the only fixed point is $0$), and the resonance density of the transit points also rules out
any recursion-internal distinction of a value (Doc.~275, epistemic limit). The recursion side is
thereby settled, not an open question; what is open is solely the existence of the external
anchor -- whether a skew-adjoint dynamics carries precisely $\chi=\pi/2$ and thereby fixes the
depth $\delta$ for an internal reason \emph{of that dynamics} (not of the recursion); if it did,
$2/9$ would acquire a dynamical,
protected privilege, otherwise it stays contingent. The elimination chain (Doc.~282) has ruled
out every static candidate -- symmetric, counting-topological, radial, static-geometric -- and
the only survivor is the dynamical skew-adjoint phase, whose origin lies outside the static
geometry. It is there, not in the magnitude and not in a flat winding number, that the value
$2/9$ would have to be sought.

\section{Placement: usefulness and proximity to application}
The findings of this document are structure and exclusion results, not technology. Their
immediate usefulness lies with three groups. First, work on the Koide relation and flavour
structure: the no-go ($2/(9\pi)$ irrational; magnitude and integer winding counter-blind)
closes off topological and counting derivation routes for $\theta=2/9$ and spares their further
pursuit; the elimination chain localizes the remaining search precisely in the dynamical
skew-adjoint phase channel. Second, methodologically: the magnitude/phase decomposition serves
as a shared audit language for framework comparisons (script-backed claims, pre-registered
criteria, matched nulls) and transfers independently of the FFGFT content. Third, didactically:
the translation between particle physics and communications engineering (circulant as DFT,
all-pass, injection locking, Arnold tongues) makes the phase channel computable with standard
engineering tools.

Only the phase channel itself touches measurement: the structural statement that mixing and
$CP$ violation live in the skew-adjoint channel (Doc.~289) has a falsifiable form in the
oscillation phases and $\delta_{CP}$, which running and upcoming experiments (DUNE,
Hyper-Kamiokande) measure. An application in the narrow sense -- a device, a procedure,
short-term engineering use -- does not follow: the tools run from communications engineering
into physics, not back, and as long as $\delta$ is not derived, the yield is that of a
sharpened open problem, not of a result.
